\notechapter{N-20230316}{Fixed Points, Iteration, and the Banach Contraction Principle}

\section{Fixed points}

A fixed point of a map \(T:X\to X\) is a point \(x\in X\) such that
\[
  T(x)=x.
\]
Intuitively, applying the transformation does not move the point.

A basic example is
\[
  T(x)=\cos x.
\]
The equation
\[
  x=\cos x
\]
has a unique real solution, approximately
\[
  x\approx0.739085\ldots,
\]
often called the Dottie number.  Iteration
\[
  x_{n+1}=\cos x_n
\]
converges to this fixed point for many starting values.

\section{Iteration}

Given a map \(T:X\to X\), an iterative process has the form
\[
  x_{n+1}=T(x_n).
\]
If \(x_n\to x\) and \(T\) is continuous, then
\[
  T(x)=T\left(\lim_{n\to\infty}x_n\right)=\lim_{n\to\infty}T(x_n)=\lim_{n\to\infty}x_{n+1}=x.
\]
Thus the limit of a convergent iteration is automatically a fixed point.

The hard question is not this final step, but why the iteration converges at all.

\section{Metric spaces and contractions}

\begin{definition}
A metric space is a set \(X\) equipped with a distance function \(d:X\times X\to\mathbb R_{\ge0}\) such that:
\begin{enumerate}[(i)]
  \item \(d(x,y)=0\) iff \(x=y\);
  \item \(d(x,y)=d(y,x)\);
  \item \(d(x,z)\le d(x,y)+d(y,z)\).
\end{enumerate}
\end{definition}

\begin{definition}
A map \(T:X\to X\) is a contraction if there exists a constant \(0<c<1\) such that
\[
  d(Tx,Ty)\le c\,d(x,y)
\]
for all \(x,y\in X\).
\end{definition}

\section{Banach fixed point theorem}

\begin{theorem}[Banach contraction principle]
Let \(X\) be a nonempty complete metric space, and let \(T:X\to X\) be a contraction.  Then \(T\) has a unique fixed point \(x^*\).  Moreover, for any starting point \(x_0\in X\), the sequence
\[
  x_{n+1}=T(x_n)
\]
converges to \(x^*\).
\end{theorem}

\begin{proof}
Let \(x_{n+1}=T(x_n)\).  Since \(T\) is a contraction,
\[
  d(x_{n+1},x_n)=d(Tx_n,Tx_{n-1})\le c\,d(x_n,x_{n-1}).
\]
Inductively,
\[
  d(x_{n+1},x_n)\le c^n d(x_1,x_0).
\]
For \(m>n\), the triangle inequality gives
\[
  d(x_m,x_n)\le \sum_{k=n}^{m-1}d(x_{k+1},x_k)
  \le \sum_{k=n}^{m-1}c^k d(x_1,x_0)
  \le \frac{c^n}{1-c}d(x_1,x_0).
\]
Thus \((x_n)\) is Cauchy.  Completeness gives a limit \(x^*\).  By continuity of \(T\),
\[
  T(x^*)=x^*.
\]
For uniqueness, if \(Tx=x\) and \(Ty=y\), then
\[
  d(x,y)=d(Tx,Ty)\le c\,d(x,y).
\]
Since \(c<1\), this forces \(d(x,y)=0\), so \(x=y\).
\end{proof}

\section{Error estimate}

The proof gives an explicit estimate:
\[
  d(x_n,x^*)\le \frac{c^n}{1-c}d(x_1,x_0).
\]
This is useful computationally because it tells us how quickly the iteration converges.

\Needspace{16\baselineskip}
\section{Applications}

\begin{example}[Solving equations]
To solve \(f(x)=0\), one may rewrite it as \(x=T(x)\).  If \(T\) is a contraction on a complete interval or metric space, iteration converges to the solution.
\end{example}

\begin{example}[Differential equations]
Picard iteration for ordinary differential equations is an application of the contraction principle in a function space.  The local existence and uniqueness theorem for ODEs can be proved this way.
\end{example}

\begin{example}[Functional equations]
Many recursive definitions can be interpreted as fixed point problems.  In computer science, fixed point theorems appear in domain theory, semantics of recursion, and program analysis.
\end{example}

\section{Beyond Banach}

Banach's theorem is powerful because it gives existence, uniqueness, and an algorithm.  Other fixed point theorems weaken some of these conclusions:
\begin{enumerate}[(i)]
  \item Brouwer's fixed point theorem gives existence for continuous maps from a closed ball to itself, but not uniqueness or iteration convergence;
  \item Schauder's fixed point theorem extends Brouwer-type ideas to infinite-dimensional settings;
  \item Tarski's fixed point theorem concerns monotone maps on complete lattices.
\end{enumerate}

\section{What the example reveals}

The fixed point idea is a bridge between analysis, topology, computation, and dynamics.  The Banach contraction principle is the cleanest analytic version: contraction gives a unique fixed point and a concrete iterative algorithm for finding it.


\section{Why completeness is the hidden hypothesis}

The contraction estimate shows that the iterative sequence is Cauchy.  It does not by itself produce a point inside the space.  Completeness is exactly the condition that Cauchy sequences have limits in the space.  This is why the theorem is false on an incomplete space: the iteration may approach a missing point.

For example, in \((0,1)\) one may construct a contraction whose natural limit is \(0\), but \(0\notin(0,1)\).  The shrinking process is not enough unless the space contains its limiting points.

\section{Picard iteration as fixed-point theory}

The same idea solves ordinary differential equations.  The initial value problem
\[
  y'(t)=f(t,y(t)),\qquad y(t_0)=y_0
\]
is equivalent to
\[
  y(t)=y_0+\int_{t_0}^{t}f(s,y(s))\,ds.
\]
Define an operator \(T\) on functions by the right-hand side.  Under a Lipschitz condition and on a short interval, \(T\) is a contraction.  Its fixed point is the solution.  Thus the theorem is not merely about numerical iteration; it is the engine behind existence and uniqueness.

\section{Source repair: non-contractive fixed-point theorems}

The uploaded note does not stop at Banach's theorem.  It asks what happens when the map is not a strict contraction.  Krasnoselskii-type theorems apply when an operator can be decomposed as
\[
  T=A+B,
\]
where \(A\) is contractive and \(B\) is compact, usually on a closed convex subset of a Banach space.  The guiding idea is that the contractive part controls distance, while the compact part supplies finite-dimensional approximation behavior.

Meir--Keeler's theorem weakens the fixed contraction constant.  Instead of requiring one number \(q<1\) for all distances, it asks that distances in a small band above \(\varepsilon\) are pushed below \(\varepsilon\).  Geraghty-type theorems further replace a fixed constant by a control function.  The common theme remains: some form of distance decrease forces Cauchy behavior.

\section{Source repair: Brouwer and Schauder}

Brouwer's theorem says that a continuous map from a closed ball in finite-dimensional Euclidean space to itself has a fixed point.  It gives existence but not uniqueness, and it does not give the Banach iteration estimate.

Schauder's theorem is an infinite-dimensional analogue: compact continuous maps on suitable closed convex sets have fixed points.  This is important in differential and integral equations, where solutions are often obtained as fixed points of compact operators.

The contrast is useful:
\[
  \text{Banach: contraction }\Rightarrow\text{ unique fixed point and iteration},
\]
\[
  \text{Brouwer/Schauder: compactness and convexity }\Rightarrow\text{ existence only}.
\]
